% Section 2, Background and Motivation: Pierazzi et al.'s problem space, constraints, reachable set, and side effects;
% who needs which description of the reachable set, and the gap they share. Spine: docs/framework-spine.md (P1-P4),
% with the user's decisions and comments of 2026-11-04 (no domain examples while the theory is drafted).

We follow the formalization of problem-space attacks by Pierazzi et al.~\cite{pierazzi2020intriguing}. A classifier
does not see the objects of its domain, such as programs, documents, or network traffic, but vectors of numbers
computed from them. The objects form the \emph{problem space} $\mathcal{Z}$, the vectors form the \emph{feature space}
$\mathcal{X} \subseteq \mathbb{R}^n$, and a \emph{feature mapping} $\varphi: \mathcal{Z} \to \mathcal{X}$ takes each
object $z$ to its feature vector $\mathbf{x} = \varphi(z)$. An evasion attack changes a malicious object's
representation until the classifier labels it benign, and it can work in either space. A \emph{feature-space attack}
perturbs the vector directly, searching for a perturbation $\boldsymbol{\delta}$ such that $\mathbf{x} +
\boldsymbol{\delta}$ is misclassified. Its search is cheap, since gradients and optimizers work on vectors, but
nothing guarantees that any object has the features $\mathbf{x} + \boldsymbol{\delta}$. A \emph{problem-space attack}
edits the object itself with a \emph{transformation} $T: \mathcal{Z} \to \mathcal{Z}$, and the classifier then sees
$\varphi(T(z))$, the features of an object that exists.

Not every transformation serves the attacker. Pierazzi et al.\ name four kinds of \emph{problem-space constraints}
$\Gamma$ that a transformation must satisfy. It must be one of the transformations available to the attacker. It must
preserve the semantics that matter, such as a program's malicious behavior. It must keep the object plausible, so
that the result does not look altered to an analyst who inspects it. And it must be robust to preprocessing, so that
the non-learned processing a detector may apply before extracting features, such as removing dead code, does not
undo it. The problem-space constraints determine the \emph{feature-space constraints} $\Omega$, the restrictions on
which perturbations of $\mathbf{x}$ correspond to valid objects. Let $\mathcal{T}$ be the set of transformations,
and of finite sequences of them, that satisfy $\Gamma$. The \emph{reachable set} of $z$ is the set of feature vectors
these transformations produce,
\[
  \mathcal{R}(z) = \{\, \varphi(T(z)) : T \in \mathcal{T} \,\} \subseteq \mathcal{X}.
\]
An evasion is \emph{realizable} if the misclassified vector lies in $\mathcal{R}(z)$, that is, if some valid
transformation of the object produces it. A feature-space attack can return vectors outside $\mathcal{R}(z)$, and
$\Omega$ is the means of keeping it inside.

A problem-space attack often starts from a feature-space target: a perturbation $\boldsymbol{\delta}^*$ that would
evade the classifier, which the attacker then tries to produce with a transformation. Pierazzi et al.\ write the result
as $\varphi(T(z)) = \mathbf{x} + \boldsymbol{\delta}^* + \boldsymbol{\eta}$, where the \emph{side-effect features}
$\boldsymbol{\eta}$ are the difference between what the transformation produced and what the attacker aimed for. Their
account of where side effects come from is the constraints. To keep the object valid, a transformation may need
\emph{consolidation}, changes beyond the intended one, such as importing the dependencies of an added API call or
declaring the classes it uses, and these changes move features of their own. Consolidation is not the only source.
The feature mapping does not follow the attacker's intent, and one changed element of the object can be read by many
features, so an edit that needs no consolidation can still move features it did not aim at. Nor are side effects
inevitable: an edit that moves exactly the features it targets has $\boldsymbol{\eta} = \mathbf{0}$. Pierazzi et al.\
observe that when the feature mapping is neither differentiable nor invertible, as is usually the case in malware
detection, the side effects of a transformation generally cannot be computed in advance, and attacks find them by
applying each candidate transformation and extracting the features of the result. Whatever their source, once a
transformation is fully specified, consolidation included, the vector it produces is fixed by the feature mapping
alone, so computing its side effects is a question about $\varphi$.

% EXAMPLE (removed, D2): appending content to a file changes every feature normalized by the file's size, although
% the content touches none of the structures those features describe (EMBER: every section's size ratio falls).

Each line of work that cares whether an evasion is realizable needs some description of the reachable set, and each
needs a different one. A problem-space attack that searches over candidate transformations needs, for each candidate,
the point $\varphi(T(z))$ it reaches, side effects included. The attacks against deployed malware detectors obtain
it by applying the candidate and extracting the features of the
result~\cite{demetrio2021functionality,song2022mab}, which costs one extraction per candidate and says nothing about
the candidates the search never tries. A feature-space attack that respects realizability needs $\Omega$ written out
as constraints over features, which domain-independent attacks then take as
input~\cite{simonetto2022unified,chernikova2022fence}. The constraints are written by hand with a domain
expert~\cite{dyrmishi2023empirical} or learned from data~\cite{sheatsley2021robustness}. A detector meant to resist
problem-space attacks needs to know, for each feature, whether any transformation in $\mathcal{T}$ can lower it, on
every object the detector will face, since a detector whose score never falls when the attacker's edits raise
features cannot be evaded by those edits~\cite{incer2018adversarially,gimenez2025certifiably}. These directions are
assigned by hand~\cite{incer2018adversarially,gimenez2025certifiably} or learned from the attacks that have been
observed~\cite{gimenez2025certifiably}. All three descriptions come from outside the feature mapping: from trial, from
an analyst's reading of the extractor, or from the attacks someone happened to run. Trial is exact but answers only for
the transformations tried. A hand analysis is as complete as its author's reading of thousands of features, and a
description learned from attacks covers only what those attacks did. Yet $\mathcal{R}(z)$ is determined by $\varphi$
once $\mathcal{T}$ is fixed, and the feature mapping is code that the defender has. In the next section we compute
two descriptions from that code, for a broad class of feature mappings: the exact point that one transformation
reaches, and, for every transformation a threat model allows, a sound bound on the direction in which each feature
can move, which is both a feature-space constraint and what a monotone detector needs.
